% Standalone section for the PA195 article. Requires amsmath, amssymb,
% amsthm, hyperref and theorem/lemma/proposition environments in the parent.
% All claims below are mathematical proofs. They are not claims of a newly
% compiled Lean endpoint for the new literal machine.

\section{The PA-provable consistency equivalence}
\label{sec:pa-calibration}

Let $M_*$ denote the delivered 195-state table, with its designated blank
initial configuration, and let $H_*$ be the arithmetic assertion that this
literal machine has a finite halting trace. Throughout this section,
$\operatorname{Con}(\mathrm{PA})$ means that a fixed standard, recursively
checked first-order PA proof system has no proof of $0=1$. Changing between
the usual PA proof codings gives primitive-recursive proof translations
whose correctness PA proves. The conclusion concerns the arithmetic
sentence for the exact table, rather than an external assertion about its
standard execution.

\begin{theorem}\label{thm:pa-biconditional}
For the table $M_*$ in this package,
\[
 \mathrm{PA}\vdash
 \bigl(\neg H_*\leftrightarrow\operatorname{Con}(\mathrm{PA})\bigr).
\]
This is an ordinary mathematical proof of a PA-provability statement.
The new endpoint has not been compiled as an additional Lean or Rocq
theorem. The executable certificates supply the finite literal-program
facts used in the proof; the arguments below explain their arithmetic
meaning and the uniform proof transformations.
\end{theorem}

There are three levels to connect: PA proofs, the numerical proof calculus
searched by the register program, and literal binary-tape traces. The
published PA196 development supplies the PA-refutation compiler for the
original calculus~\cite{chang}. Its other internal direction is not
assumed here. We first justify the incoming diagonal replacement used
by $M_*$~\cite{diagonal}, and then supply the reverse proof
transformation and the finite-trace argument, developing the syntactic
route identified in the earlier bridge notes~\cite{bridge}.

\subsection{The numerical calculus and its diagonal replacement}

Write
\[
 p(a,b)=\frac{(a+b)(a+b+1)}2+b,
 \quad I(A,B)=4p(A,B),\quad M(a,b)=4p(a,b)+1,
\]
\[
 U(a,B)=4p(B,a)+2,\qquad E(a,b)=4p(a,b)+3.
\]
The notational convention is $\neg A=I(A,1)$, conjunction and existential
quantification being their classical abbreviations. Code $0$ is truth,
and code $1=M(0,0)$ is the designated contradiction. The quantifier
constructor has the transposed Cantor arguments displayed above.

The diagonal calculus $\mathcal C_\Gamma$ retains the six PA196 heads,
the dominated propositional head, modus ponens, and closed
generalization. Its guarded gate is
\begin{align*}
 \Gamma(B)&=I(B,U(B,B)),\\
 V(a,b,c)&=I(M(a,b),U(c,M(a,b))),\\
 P(a,b,c)&=I(c,M(c,b)),\\
 S(a,b,c)&=U\bigl(b,
       (M(a,b)\to P(a,b,c))\wedge(M(c,b)\to c)\bigr),\\
 Q(a,b,c)&=U\bigl(c,S(a,b,c)\to\neg P(a,b,c)\bigr),\\
 F_\Gamma(a,b,c)&=\neg\bigl(\Gamma(c)\to(V(a,b,c)\to Q(a,b,c))\bigr).
\end{align*}
The gate is available when $c\ne a,b$; a failed guard emits $0$.
The logical rules and remaining axiom families are exactly those of
\texttt{PA196/Logic/Basic.lean}, with this one replacement of its
\texttt{bundle}.

Every recursive formula child of a nonzero code is smaller than that
code. Every variable label, including a bound label, occurring in a
formula of code $B$ is strictly below $B$. These assertions follow from
the inequalities $a,b\le p(a,b)$ and the four tags, with code zero handled
separately. Consequently $x_B$ is absent from the decoded formula $B$.
The diagonal clause is a valid vacuity instance, although its binder is
chosen from the numerical code of the entire formula.

Classical implicational reasoning extracts from the gate both
$\Gamma(c)$ and the old gate $G(a,b,c)=\neg(V(a,b,c)\to Q(a,b,c))$.
Every diagonal instance $\Gamma(B)$ is available: use $c=B$ and
$a=b=B+1$. Thus all membership-vacuity instances and pointed separation
used in the original proof remain available.

\begin{lemma}[Recovery of the required equality interfaces]
\label{lem:pa-diagonal-interfaces}
The diagonal calculus derives, uniformly in the numerical labels,
\begin{align*}
 E(a,b)&\to\bigl(M(d,a)\leftrightarrow M(d,b)\bigr),\\
 E(a,b)&\to\bigl(M(a,d)\leftrightarrow M(b,d)\bigr),\\
 E(a,b)&\to R(a,b,z)\quad (z\ne a,b),
\end{align*}
where
\[
 R(a,b,z)=U\bigl(z,M(a,z)\leftrightarrow M(b,z)\bigr)
 \wedge U\bigl(z,M(z,a)\leftrightarrow M(z,b)\bigr),
\]
and it derives $\exists x_a R(a,b,z)$ for $z\ne a,b$.
All these derivations are supplied by primitive-recursive proof templates,
with correctness provable in PA.
\end{lemma}

\begin{proof}
The two forward membership rewrites are retained heads. To obtain the
reverse directions, first assume $a\ne b$. Take
$\chi(t)=M(d,t)$ or $M(t,d)$, initially with $d\ne a$, or take the
diagonal family $\chi(t)=M(t,t)$. Put
\[
 C=E(a,b)\to(\chi(b)\to\chi(a)),\qquad
 \varphi=\neg C,\qquad c=\ulcorner\varphi\urcorner.
\]
Here $c$ is the actual numerical code of the displayed formula, so it is
larger than every label in $C$. In particular, all required guards hold.
Apply the old gate at $(a,a,c)$ and abbreviate its $S,P,Q$ by those
letters. We need only its consequence $\neg Q$, not an unproved general
rule for changing bound-variable names.

The second conjunct in $S$ gives
$S\to U(a,M(c,a)\to\varphi)$. Propositionally,
$\varphi$ implies $E(a,b)$, $\chi(b)$, and $\neg\chi(a)$; hence universal
monotonicity gives
\[
 S\to U(a,M(c,a)\to\neg\chi(a)).
\]
The equality atom has disappeared from this quantified consequent.
The unchanged negative-property specialization lemma now gives
\[
 S\to\bigl(M(c,b)\to\neg\chi(b)\bigr).                 \tag{1}
\]
For precision, this specialization uses only the retained witness
$\exists x_a E(b,a)$, direct rewriting under $E(b,a)$, and vacuity of
$M(c,b)\to\neg\chi(b)$ at $a$. The latter follows from
membership vacuity because $a\ne b,c$ and, in the two nondiagonal
families, $a\ne d$. It is exactly the construction
\texttt{negativeSpecialization} in \texttt{PA196/Logic/Reverse.lean};
none of its premises is primitive equality vacuity.

Under $\varphi$, the point condition $P$ gives $M(c,a)$.
Direct element-slot rewriting under $E(a,b)$ then gives $M(c,b)$.
Together with (1), $S$ contradicts the consequence $\chi(b)$ of
$\varphi$. Classical deduction gives
$\varphi\to(S\to\neg P)$. Generalize this theorem at $c$, distribute
the quantifier through implication, and use
$\Gamma(\varphi)=\varphi\to U(c,\varphi)$. The result is
$\varphi\to Q$. Since the gate supplies $\neg Q$, double-negation
elimination gives $C$, as required.

If $a=b$, the reverse rewrite is propositional. When the fixed label
$d=a$, direct rewriting first takes $M(a,b)$ or $M(b,a)$ to $M(b,b)$;
the already proved diagonal case takes $M(b,b)$ to $M(a,a)$.
This covers every label coincidence without a freshness assumption on
the theorem being derived.

For the neighborhood formula choose $z_0=\ulcorner E(a,b)\urcorner$.
The available diagonal clause $E(a,b)\to U(z_0,E(a,b))$ and the two
bidirectional rewrites yield $E(a,b)\to R(a,b,z_0)$.
To change the binder to an arbitrary $z\ne a,b$, only the case
$z\ne z_0$ remains. Let $P(t)$ be either membership-only matrix in
$R$. Bidirectional rewriting gives
$E(z,z_0)\to(P(z_0)\to P(z))$.
The witness $\exists x_{z_0}E(z,z_0)$, universal distribution, and
vacuity of $P(z)$ at $z_0$ give
\[
 U(z_0,P(z_0))\to P(z).
\]
Its antecedent is vacuous at $z$: its atoms contain only
$a,b,z_0$, and membership vacuity extends through its Boolean
connectives and the outer quantifier by the retained commutation head.
Generalization and distribution therefore give
$U(z_0,P(z_0))\to U(z,P(z))$.
This proves the claimed neighborhood implication for every allowed $z$.
Finally the unchanged witness head and existential monotonicity give
the neighborhood witness, reversing its two arguments by the
propositional symmetry of $R$ if necessary.

Every operation in this proof is a fixed syntactic proof template,
composition, numerical code formation, or a displayed decidable
label-equality case split. The choices $c=\ulcorner\neg C\urcorner$ and
$z_0=\ulcorner E(a,b)\urcorner$ use no unbounded search. PA verifies the
templates by the same proof-checker induction as their original
counterparts. This establishes the asserted uniformity.
\end{proof}

The PA196 refutation compiler uses primitive equality through these
membership-rewriting, neighborhood, and witness interfaces. In its
logical and arithmetic certificate copies, the uses of equality vacuity
are in the reverse-rewriting and neighborhood modules just replaced.
The old gate, all other heads, and the command stack interface remain
available. Substituting Lemma~\ref{lem:pa-diagonal-interfaces}'s templates
therefore gives a PA-verifiable refutation compiler for
$\mathcal C_\Gamma$. In particular,
\begin{equation}\label{eq:pa-forward-calculus}
 \mathrm{PA}\vdash
 \exists p\,\operatorname{Prf}_{\mathrm{PA}}(p,\ulcorner0=1\urcorner)
 \ \longrightarrow\
 \exists d\,\operatorname{Prf}_{\Gamma}(d,1).
\end{equation}
This is reuse of the published proof-producing construction with explicit
replacement templates, rather than an appeal to standard-model
equiconsistency. The 191-item AX word still pushes nine items in the
same positions as the original nine-item command, so its stack
compiler indices need no alteration.

\subsection{A uniform translation of the numerical calculus into PA}

We now prove the converse proof-theoretic direction. Work first in the
fixed conservative definitional extension of PA containing
$\operatorname{pow2}(n)=2^n$ and its usual arithmetical definition.
Only fixed primitive-recursive functions are added. A fixed
primitive-recursive elimination of these definitions converts the
resulting proofs to ordinary PA proofs.

Use the bit predicate
\[
 B(s,z)\;\equiv\;
 \exists u,v\,[s=(2u+1)2^z+v\ \wedge\ v<2^z].
\]
PA proves the usual digit lemmas directly by induction and division
with remainder. In particular $B(s,z)\to z<s$, since
$s\ge 2^z\ge z+1$, and hence PA proves $\neg B(x,x)$.

Let $T(c)$ be the arithmetic formula obtained by decoding the source
code $c$:
\begin{align*}
 T(0)&=(0=0),& T(M(a,b))&=B(x_a,x_b),\\
 T(E(a,b))&=(x_a=x_b),& T(U(a,A))&=\forall x_a\,T(A).
\end{align*}
For a nonzero implication code the translation is the corresponding
implication. There is one numerical alias:
$I(0,0)=0$. Thus the uniform constructor theorem is a proof of
\[
 T(I(A,B))\leftrightarrow(T(A)\to T(B)),
\]
not an unconditional assertion of syntactic identity. At $A=B=0$ it
uses the fixed proof $(0=0)\leftrightarrow((0=0)\to(0=0))$.
Every other case follows from the definition of $T$.

Arithmetic auxiliary bound variables are represented with de Bruijn
indices, or equivalently renamed above all source labels before
substitution. The source label $c$ is therefore absent among the free
variables of $T(c)$ even though the expanded bit formula has its own
bound variables. Quantifier formation closes the indicated free source
label, so nested occurrences of the same label have their usual
shadowing semantics. No assumption that all numerical labels are
distinct is made.

\begin{lemma}[Uniform finite bit comprehension]
\label{lem:pa-bit-comprehension}
Given the syntax code of any arithmetic formula $\theta(z,\vec u)$,
there is a primitive-recursive procedure producing a PA proof of
\[
 \forall n\,\exists s<2^n\,\forall z<n\,
       \bigl(B(s,z)\leftrightarrow\theta(z,\vec u)\bigr).
                                                        \tag{2}
\]
PA proves that the procedure produces a valid proof for every
well-formed formula input, with the indicated fresh variables.
\end{lemma}

\begin{proof}
Take induction on $n$ with the existential formula in (2) as induction
formula. At zero choose $s=0$. Given a witness $s<2^n$, split by
excluded middle on $\theta(n,\vec u)$. If it holds use $s+2^n$;
otherwise use $s$. The fixed bit lemmas say that the first $n$ bits
are unchanged, that the new $n$th bit is respectively one or zero,
and that the result is below $2^{n+1}$. This is a fixed finite
first-order proof pattern with one PA induction instance.

On syntax codes the procedure performs capture-avoiding substitution,
forms the displayed induction instance, and concatenates this fixed
proof pattern with the fixed arithmetic lemma proofs. It never
evaluates $\theta$. Its syntax operations and a fresh-bound-variable
choice are primitive recursive. The PA proof of correctness is
induction on the syntax operations, followed by direct verification
that the generated induction line is an instance of the PA induction
schema and that each subsequent line follows by the stated logical
rule. Thus this is a statement about proof checking, with an arbitrary
formula code as input, not a uniform satisfaction or reflection
principle for PA. Full PA induction is used: $\theta$ is not restricted
to a bounded or quantifier-free formula.
\end{proof}

We spell out the gate supplier because label collisions matter here.
Fix $a,b,c$ with $c\ne a,b$, let $\phi=T(c)$, and regard $x_a,x_b$
as the free values before the inner quantifier is entered.
If $a\ne b$, apply (2) with bound $x_a+x_b+1$ and matrix
\[
 \bigl(B(x_a,z)\vee z=x_b\bigr)\wedge\phi[z/x_b].
                                                        \tag{3}
\]
The witness $s$ satisfies, for every $z$,
\[
 B(x_a,z)\to(\phi[z/x_b]\to B(s,z)),\qquad
 B(s,z)\to\phi[z/x_b],
\]
and $\phi\to B(s,x_b)$. To pass from the bounded comprehension
formula to the displayed unbounded assertions, use
$B(x_a,z)\to z<x_a$ and the fact that a number below $2^n$ has no
set bit at an index at least $n$.

If the \emph{labels} $a=b$ coincide, the first antecedent inside
$T(S)$ becomes $B(z,z)$, not $B(x_a,z)$ with $x_a$ held fixed.
It is therefore false. Apply (2) instead with bound $x_b+1$ and
matrix
\[
 (z=x_b)\wedge\phi[z/x_b].                              \tag{4}
\]
This witness has the required second implication and distinguished
point condition, and the first implication is vacuous. Equality
substitution in (4) supplies $\phi$ at the distinguished point.
The substitutions in (3)--(4) are capture avoiding; in particular
the inner binder corresponding to $b$ is renamed when the witness
expression depends on the original free value of $x_b$.

In both cases the result is a PA proof of
$\exists x_c\,(T(S)\wedge T(P))$. Freshness of $c$ in $\phi$
ensures that choosing the set witness does not alter the property.
Consequently PA proves $\neg T(Q)$. The guarded membership-vacuity
clause has a first-order proof, and $T(\Gamma(c))$ has a first-order
proof by the code-freshness lemma. Classical propositional reasoning
therefore proves $T(F_\Gamma(a,b,c))$.
When the guard fails the command emits $0$, with its fixed truth proof.

The source false code is translated to $T(1)=B(x_0,x_0)$.
Its negation is a fixed PA theorem, uniformly under all binders.
Thus source negation, conjunction, and existential abbreviations have
PA-provable equivalences to the ordinary classical connectives, even
if a surrounding source quantifier binds label zero. This deals with
the false-code convention independently of the truth-code alias.

The other head suppliers are fixed first-order proof templates:
the witness-elimination head, quantifier commutation, the two equality
congruence heads, an equality witness, and quantifier distribution.
Coincident labels are handled by capture-avoiding instantiation, with
a numerical equality case split if needed. The dominated
propositional head is a classical tautology; its arithmetic validity
does not require the dominance restriction. Generalization and modus
ponens use their ordinary PA proof constructors, together with the
implication-constructor equivalence above.

\begin{lemma}[Syntactic refutation translation]
\label{lem:pa-refutation-translation}
There is a primitive-recursive function $J$ such that
\[
 \mathrm{PA}\vdash\forall d\,
 \left(\operatorname{Prf}_{\Gamma}(d,1)
 \to\operatorname{Prf}_{\mathrm{PA}}
             (J(d),\ulcorner0=1\urcorner)\right).
                                                        \tag{5}
\]
\end{lemma}

\begin{proof}
Decode the source derivation by recursion on its finite proof code.
At an axiom node insert the appropriate supplier just described.
At modus ponens concatenate the two translated proofs and the
compositional implication proof; at generalization append the
corresponding quantifier rule. At a code-zero alias emit or use the
fixed truth equivalence. Recursion on a malformed proof is made total
by an arbitrary fixed output; the correctness statement has the
displayed proof-validity hypothesis.

PA proves the syntactic invariant
$\operatorname{Prf}_{\Gamma}(d,A)\to
  \operatorname{Prf}_{\mathrm{PA}}(J_0(d),T(A))$
by induction on this coded recursion. Constructor children decrease
strictly; proof concatenation, label closing and fresh-variable
renaming have primitive-recursive definitions and their usual
syntactic correctness inductions. The gate case is precisely
Lemma~\ref{lem:pa-bit-comprehension} and the two label cases above.
Append the fixed theorem $\neg B(x_0,x_0)$ at a proof of code one,
then eliminate the fixed definitional extension. This gives (5).

This argument proves one quantified assertion about proof codes
inside PA. It does not merely produce a PA proof separately for each
standard formula and then assert a uniform conclusion. The inductive
verification of the compiler, including the schema supplier, is the
step that makes (5) applicable to nonstandard codes in every model
of PA.
\end{proof}

\subsection{The exact proof-store invariant}

The numerical command store is $(t,w,d,h)$. Both $t$ and $h$ carry
formula codes; $w$ is the Cantor-coded stack and $d$ is an arithmetic
parameter. Annotate $t,h$ with PA proofs of $T(t),T(h)$, and annotate
each live stack element with the corresponding proof. The clean store
uses the fixed proof of $T(0)$.

Here is a precise convention that avoids treating Cantor pairing as a
free list constructor. Let $\operatorname{hd},\operatorname{tl}$ be
the two Cantor projections. At zero both are zero. For $w>0$,
$\operatorname{tl}(w)<w$: writing $w=p(a,b)$, if $b>0$ then
$p(a,b)>b$, and if $b=0$ the inequality is immediate. Define
\[
 \delta(w)=\min\{k\le w:\operatorname{tl}^{k}(w)=0\}.
\]
PA proves the set nonempty by induction on $w$.
The canonical stack has exactly these $\delta(w)$ heads, and its
proof annotation is a sequence of that length. Its invariant is
\[
 \forall i<\delta(w)\quad
 \operatorname{Prf}_{\mathrm{PA}}
 \bigl(\pi_i,T(\operatorname{hd}(\operatorname{tl}^{i}(w)))\bigr).
                                                        \tag{6}
\]
Bounded minimization and bounded sequence operations make this an
ordinary primitive-recursive invariant.

Pushing a code $a$ onto $w$ prepends its proof, except when $a=w=0$;
in that one case $p(0,0)=0$ and the canonical annotation remains
empty. Popping nonzero $w$ removes the first proof. Popping zero
returns code zero with its fixed truth proof and leaves an empty
annotation. These statements are exact consequences of
$\operatorname{decode}(p(a,b))=(a,b)$ and
$p(a,b)=0\leftrightarrow a=b=0$, and PA proves each of them.

Now inspect the actual command updates.
AX pushes the saved top, a generalization of the old $h$, and the
seven ordinary head instances, in their specified order, and leaves
the guarded diagonal gate in $t$ with $d=h=0$.
The nine pushed entries have the PA proofs provided above.
For MP with offset $x$, put $y=d+x$.
If $I(t,y)=h$, combine the proofs for $t,h$ to obtain a proof for
$y$, push the truth code, and set $(t,d,h)=(y,y,0)$.
If $I(t,y)=h=0$, then $t=y=0$ and the same operation is justified
by the truth proof. On a failed match the command outputs truth,
retains the stack, and sets $(t,d,h)=(0,y,0)$.
Shift, when it does not halt, sets
\[
 (t,w,d,h)\longmapsto
   \bigl(h,\operatorname{tl}(w),d,\operatorname{hd}(w)\bigr),
\]
and moves the annotations accordingly.

The literal Shift test is $t+h=1$. PA proves that this implies
$t=1$ or $h=1$. In either case the relevant annotation, followed by
the fixed proof $\neg T(1)$, yields a PA refutation. In particular
the invariant must track both visible formula registers, rather
than just $t$.

Every annotation update is primitive recursive. The local statement
just proved holds for \emph{every annotated store}: a successful
command preserves the annotation invariant, and a halting command
produces a PA refutation. PA induction therefore propagates the
invariant through an arbitrary coded finite command sequence,
including one starting at an already live scheduler store.
In particular, if execution from the clean store halts then a PA
refutation exists. Under $\operatorname{Con}(\mathrm{PA})$, every
annotated command has a nonhalting output and its successor store
is annotated again. This is the internal replacement for the
upstream external HF-model soundness argument.

\subsection{Why the finite runtime certificates are sufficient in PA}

The source and physical verifications in this package have the
following mathematical content. A frame consists of its source
path, a gap, and the complete register file. A value $v$ is stored
as $0\,1^{v+1}$. The physical gap is represented by an affine run
length $n+c$, with arbitrary $n\ge0$; register contents and the
complete tail are left symbolic. The initial file has at least ten
blocks, and this lower bound is preserved by every primitive:
increment appends one extra block and decrement retains the length.
The greatest source path length is 37, below the header constant 48.

For each increment, each jump, each zero and positive decrement
branch, and the halt site, the delivered finite trace certificate
starts at the exact frame and ends at the exact successor frame
or at the distinguished halt control. There are 4,815 cases for
the 3,700-site source. These are symbolic branch contracts, rather
than 4,815 examples of numerical register inputs.

There are only four kinds of step in their symbolic interpreter:
an actual table row; a change of notation for an identical tape;
a whole run of one self-loop; and a register-file excursion.
Notation changes preserve every concrete instance. A self-loop
over a run of length $an+c$ has the stated effect by induction on
that length. The register excursions are marking a register,
incrementing it, zero decrement, and positive decrement.
Their proofs in \texttt{M6/Native.lean} are inductions on the finite
file and its unary blocks. They use exactly the specified engine
rows and the append-one tail row. These rows are unchanged by the
new decoder and the seven simultaneous folds.

Consequently the soundness proof of
\texttt{M6/Runtime.lean}, specifically \texttt{sstep\_sound} and
\texttt{exec\_sound}, applies to each case. It gives exact
finite reachability for every natural gap and every valid full
register file. The checker only accepts after it has traversed at
least one actual row, so each completed primitive segment has
strictly positive physical length. A nonhalting contract ends at
the root control, not at the halt control.

The same assertion is provable arithmetically, for the following
specific reason. Finite tapes are coded finite words; runs,
concatenation, frame construction, and the register operations are
primitive recursive. The proofs for a self-loop and each of the
four register excursions use induction on a natural run length or
a coded finite list. Their witness traces can be produced by the
same bounded scans and primitive recursions. Repeating the four
symbolic-step cases and concatenating their traces is again a
primitive recursion. PA proves these inductions directly. The
candidate-specific part is a finite list of table equalities and
completed symbolic traces, all decidable closed arithmetic facts.
Thus PA proves the exact, positive source-to-physical segment
assertion uniformly in coded frames. An external theorem about
standard tapes alone would not suffice; the just-described coded
induction is the internalization used here.

The AX source binding checks every primitive and relative jump in
each of its 191 macro blocks against the inherited total arithmetic
macro, and checks every exterior successor. The free-algebra emitter
check supplies its exact numerical output formulas. Hence the
source-level command and scheduler contracts are inherited with
the gate replaced as stated. The scheduler's other 169 sites are
unchanged. Its finite command-word decoding, the Cantor-pairing
routines, and the fixed AX macro composition are primitive-recursive
computations with PA-provable terminating bounds. More specifically,
the construction \texttt{parsed\_command\_step} in
\texttt{PA196/Source/CommandTransfer.lean} takes the next decoded
command at a fetch boundary, executes it, and returns to the next
fetch boundary in a strictly positive number of primitive source
steps whenever the decoded command has a nonhalting output. It applies
to every store and every scheduler pair $(n,l)$ satisfying that explicit
output premise; the annotated-store invariant under consistency
supplies the premise used below.
The source binding retains this construction with the new AX
composition. The store need not be reset between decoded words;
the proof-store invariant is preserved from one command boundary
to the next.

The complete blank startup is checked for exactly 58,154 steps,
including the entire final nonzero support, the head position,
and the state. It ends at the stated clean source frame.
In native labels that state is 84 and the head coordinate is $-2$;
the injective dense relabeling is checked separately. As a fixed
finite computation, the startup equality and the absence of an
earlier halt are PA-provable.

\begin{lemma}[Reflection of a finite physical halt]
\label{lem:pa-positive-reflection}
PA proves that a finite halting trace of $M_*$ yields a PA
refutation.
\end{lemma}

\begin{proof}
Work in PA and assume $\operatorname{Con}(\mathrm{PA})$.
At a command boundary, decode the next command by the fixed
primitive-recursive function
$\operatorname{parse}(\operatorname{seed}(n,l))$.
The proof-store invariant makes its Shift halt test impossible,
and its nonhalting update supplies the next annotated store.
The positive \texttt{parsed\_command\_step} contract just described
then supplies a finite source execution to the next command
boundary. This is exactly the strengthened syntactic version of
the \texttt{SoundBoundary}/\texttt{boundary\_progress} argument in
\texttt{PA196/Source/Nonhalting.lean}; the external appeal to an
HF model is replaced by the PA proof annotations and consistency.

For any given $T$, primitive recursion through $T+1$ such command
boundaries therefore constructs more than $T$ primitive source
steps from the clean initial source frame. Every command segment
has positive length and ends at the next nonhalt fetch boundary.
The intermediate macro routines terminate by their stated
arithmetic contracts. Thus no source halt can occur inside an
unfinished final macro either. Concatenating the exact positive
physical segments for these primitive steps constructs a physical
run from the same initial frame longer than $T$.

Now suppose there were a blank-tape halting trace of length $T$.
If it stops within startup, the fixed startup computation refutes
it. Otherwise concatenate startup and the physical segments just
constructed. Their total length exceeds $T$. Determinism forces
the alleged halted configuration to be the same configuration
at that index in the constructed run. With the partial operational
semantics there is no outgoing row at halt, a contradiction.
With an absorbing totalization of halt the argument is equivalent:
the run could never again reach a nonhalt root frame.

No inverse decoding of an arbitrary intermediate tape is used.
Nor do we need an unbounded least-time search for the next frame.
The chosen number $T+1$ of completed command segments is a
bounded coded iteration. PA proves sequence concatenation and the
determinism comparison by induction on finite trace length.
We have therefore proved in PA that consistency excludes $H_*$,
or equivalently $H_*\to\neg\operatorname{Con}(\mathrm{PA})$.
\end{proof}

For clarity, exact reachability to a nonhalt endpoint already rules
out a premature halt inside a primitive segment: a deterministic
partial run cannot pass through a state with no outgoing row.
The required runtime premise is therefore the exact positive
endpoint contract verified here, not a disjunction allowing an
unexplained early halt instead of the next source frame.

\subsection{Completing the other direction}

By (\ref{eq:pa-forward-calculus}), a PA refutation produces a
$\mathcal C_\Gamma$ refutation. The unchanged nine-item command
compiler turns that derivation into a finite command word which
halts, using the new gate supplier where required. Here one uses
the existing proof compiler's protected-stack-tail and scheduling
identities, which permit the actual live store; no reset between
candidate words is being assumed. The unchanged scheduler reaches
the compiled word, or halts earlier, by the same bounded candidate
induction as in the original compiler. The total arithmetic macro implementations and
the source-to-physical contracts then give a finite blank-tape
halt of the exact delivered table. These are proof-producing
finite computations and their correctness inductions take place
in PA. Therefore
\[
 \mathrm{PA}\vdash
 \neg\operatorname{Con}(\mathrm{PA})\to H_*.
\]
Combining this with Lemma~\ref{lem:pa-positive-reflection} and
classical contraposition proves Theorem~\ref{thm:pa-biconditional}.

\subsection{Logical strength and the role of ConciseIPRA}

The proof does not replace PA by intuitionistic primitive recursive
arithmetic, the subject of the ConciseIPRA project~\cite{ipra}.
The finite-comprehension supplier uses an induction
instance for an arbitrary formula, which is precisely why a short
iPRA proof checker cannot simply be substituted as the searched
theory. The transferable ideas are smaller derivable interfaces,
freshness read from syntax, and uniform proof-certificate
manipulation. Here the diagonal code-freshness device is already
present in the incoming PA research~\cite{diagonal}; the new reduction is the
compatible physical allocation of that source, not a new claim
that iPRA has PA's consistency strength.
